\chapter{动力学与决策窗口}
\label{ch:dynamics}

\section{意志能量}

\begin{definition}[意志能量]
意志能量 $\Will(t)$ 是指在时刻 $t$ 可用于打断当前行为框架、并朝目标框架发起稳定化行动的有限能力。
\end{definition}

\begin{definition}[杠杆作用]
自由意志在时刻 $t$ 的有效杠杆作用定义为
\[
  \Leverage(t) = \frac{\Will(t)}{\DV(t)}.
\]
\end{definition}

\begin{theorem}[决策窗口原理]
自由意志的战略效果最大，并不发生在用力最猛烈的时候，而是发生在激活障碍相对于可用意志能量达到局部最小时。
\end{theorem}

\begin{proof}[证明思路]
如果用 $\Will(t)/\DV(t)$ 来表示杠杆作用，那么当 $\Will(t)$ 上升或 $\DV(t)$ 下降时，杠杆作用都会上升。在日常生活中，大幅提升可用意志往往并不可靠，但障碍的暂时下降却很常见：刚醒来之后、刚完成一项任务之后、进入一个新环境之后，或在经历一次明显的情绪更新之后。因此，最可靠的策略不是等待英雄式的动力爆发，而是在障碍低谷中及时行动。
\end{proof}

\section{决策窗口}

\begin{definition}[决策窗口]
决策窗口是一个短暂区间：在这段时间里，当前状态已经失去了一部分稳定力量，而下一个状态还没有被彻底排斥。
\end{definition}

常见的窗口包括：
\begin{itemize}[nosep]
  \item 刚醒来、干扰尚未加载之前；
  \item 一个既定任务块结束后的那一分钟；
  \item 身体刚到达某个任务专属地点时；
  \item 通过社会承诺强行开启新框架的时刻。
\end{itemize}

\section{失败模式}

\begin{lemma}[延迟响应失败]
如果行动者等到某个窗口关闭、惯性重新完全建立之后才行动，那么同一个目标行动所需的努力一定会严格高于窗口仍然打开时的水平。
\end{lemma}

\begin{discussion}
这个引理解释了一种很常见的误解。人们常以为自己已经“决定”行动了，因为他们口头上认可了一个未来动作。但如果他们没有在决策窗口内部真正移动起来，状态很快又会重新变贵，之前那句口头决定在机制上几乎没有价值。
\end{discussion}

\section{窗口的比较图谱}

\begin{center}
\begin{tabularx}{\textwidth}{@{}l l X@{}}
\toprule
窗口类型 & 典型时长 & 策略用法 \\
\midrule
晨间重置 & 5--20 分钟 & 在默认输入涌入之前，先注入一份预写好的计划 \\
任务边界 & 1--5 分钟 & 在切换成本已经开始支付时，直接切入下一个已准备好的任务 \\
空间转场 & 不定 & 把目标行为绑定到那些会削弱旧线索的地点上 \\
社会触发 & 不定 & 把个人意图转化为一个已经排定的外部义务 \\
\bottomrule
\end{tabularx}
\end{center}

\section{推荐阅读}

\begin{readingbox}
\textbf{基础}
\begin{itemize}[nosep]
  \item Steven Strogatz, \emph{Nonlinear Dynamics and Chaos}.
  \item Eugene Izhikevich, \emph{Dynamical Systems in Neuroscience}.
\end{itemize}
\textbf{现代研究}
\begin{itemize}[nosep]
  \item Nature Neuroscience 中关于入睡过程中可预测分岔动力学的研究。
  \item 关于脑动力学中 metastability 的综述工作。
\end{itemize}
\textbf{前沿}
\begin{itemize}[nosep]
  \item 只有在基本的决策窗口逻辑已经清楚之后，才适合使用前沿网络模型。
\end{itemize}
\end{readingbox}

\begin{summarybox}
第二章解释了为什么单纯强调用力并不是一个好的操作原则。时机会压倒强度。这个模型真正的价值，在于定位那些可重复出现的低障碍窗口，并把它们当作稀缺资产来经营。
\end{summarybox}
